\documentclass[12pt, a4paper]{article}
\usepackage[T2A]{fontenc}
\usepackage[utf8]{inputenc}
\usepackage[russian]{babel}
\usepackage{amsmath, amsfonts, amssymb}
\usepackage{geometry}
\geometry{left=30mm, right=15mm, top=20mm, bottom=20mm}
\usepackage{graphicx}
\usepackage{listings}
\usepackage{algorithm}
\usepackage{algpseudocode}
\usepackage{tocloft}
\usepackage{indentfirst}
\usepackage{setspace}
\usepackage{booktabs}
\usepackage{tikz}
\usetikzlibrary{shapes,arrows,positioning}
\usepackage{hyperref}

\onehalfspacing

\hypersetup{
    colorlinks=true,
    linkcolor=blue,
    filecolor=magenta,
    urlcolor=cyan,
    pdftitle={Курсовая работа: BbValidator (Подготовка данных)},
    pdfpagemode=FullScreen,
}

\renewcommand{\cftsecleader}{\cftdotfill{\cftdotsep}}

% Настройка lstlisting для псевдокода
\lstset{
    basicstyle=\small\ttfamily,
    breaklines=true,
    frame=tb,
    numbers=left,
    numberstyle=\tiny,
    keywordstyle=\color{blue}\bfseries,
    commentstyle=\color{gray}\itshape,
    stringstyle=\color{green!60!black},
    language=Python,
    showstringspaces=false
}

\begin{document}

\newpage
\hypersetup{pageanchor=true}

% --- АННОТАЦИЯ ---
\section*{Аннотация}
Данная курсовая работа посвящена разработке и программной реализации методов биофизического анализа трехмерной структуры остова белков (backbone) и генерации сбалансированных обучающих выборок для валидации белковых структур. Быстрый прогресс в области генерации белков (модели RFDiffusion, Chroma) привел к необходимости отсеивать физически некорректные генерации до этапа химического синтеза. Разработанный пре\-процес\-синг-ин\-тер\-фейс принимает на вход пространственные координаты атомов остова ($N$, $C_\alpha$, $C$) и рассчитывает $31$ признак узлов (двугранные углы, выбросы Рама\-чанд\-рана (Ramachandran outliers), плотность упаковки, burial ядра белка, локальную кривизну цепи, фрагментную схожесть) и $20$ признаков пар. Также в работе предложены и исследованы $7$ деформационных стратегий негативного сэмплирования (генерации декоев), имитирующих различные геометрические и топологические дефекты. Для предотвращения утечки данных (data leakage) разработано разделение данных, сгруппированное по родительским цепям исходных белков (PDB ID + Chain ID).

\newpage

% --- СОДЕРЖАНИЕ ---
\tableofcontents
\newpage

% --- ВВЕДЕНИЕ ---
\section{Введение}

\subsection{Мотивация работы}
Создание новых белков с заданными функциями (de novo protein design) является одной из ключевых задач современной синтетической биологии и медицины. Появление генеративных моделей на основе диффузии и авторегрессионного сэмплирования позволило исследователям генерировать миллионы пространственных структур белкового остова за считанные часы. Тем не менее, значительная часть генерируемого остова содержит некорректную геометрию, стерические столкновения, разорванные пептидные связи или неупакованные гидрофобные ядра.

Поскольку влажное лабораторное тестирование синтезированных белков требует колоссальных временных и финансовых ресурсов, возникает критическая необходимость в высокопроизводительных вычислительных фильтрах. Эффективность работы таких фильтров напрямую зависит от качества признаков (feature engineering) и репрезентативности обучающих наборов данных. Использование исключительно сырых декартовых координат малоэффективно из-за трансляционной и ротационной инвариантности пространственной укладки. Таким образом, разработка ротационно-инвариантного препроцессинга признаков на основе физико-химических и геометрических дескрипторов остова, а также создание сбалансированной обучающей выборки с многоуровневым негативным сэмплированием представляют собой важную практическию задачу.

\subsection{Постановка задачи}
Целью работы является создание ротационно-инвариантных методов геометрического препроцессинга белковых структур и разработка конвейера генерации обучающих выборок.
Входными данными являются пространственные координаты атомов остова $X \in \mathbb{R}^{L \times 3 \times 3}$, где $L$ — длина белковой цепи, а размерность $3 \times 3$ соответствует пространственным координатам $(x, y, z)$ трех ключевых атомов каждого остатка: азота пептидной группы ($N$), альфа-углерода ($C_\alpha$) и карбонильного углерода ($C$).

Для успешного решения задачи валидации необходимо разработать алгоритмы расчета:
\begin{enumerate}
    \item \textbf{Локальных геометрических признаков}: торсионных углов $\phi, \psi, \omega$, индикаторов конформационных выбросов Рамачандрана и координат виртуальных атомов $C_\beta$ и $O$.
    \item \textbf{Макроскопических прокси-признаков свертываемости}: плотности упаковки ядра на различных расстояниях, относительной глубины залегания остатков и локального изгиба цепи.
    \item \textbf{Попарных признаков взаимодействия остатков} и формирования разреженного $k$-ближайших соседей (kNN) графа контактов.
    \item \textbf{Синтетических дефектных структур (декоев)} разного уровня выраженности для обучения бинарных классификаторов и многоклассовых детекторов ошибок.
\end{enumerate}

\newpage

% --- ГЕОМЕТРИЧЕСКИЕ И БИОФИЗИЧЕСКИЕ ДЕCКРИПТОРЫ БЕЛКОВ ---
\section{Геометрические и биофизические дескрипторы белков}

\subsection{Торсионные углы остова}
Белки представляют собой длинные цепочки аминокислот. Поскольку длины связей и валентные углы между соседними атомами жестко зафиксированы законами химии, трехмерная форма белка однозначно определяется углами поворота вокруг связей остова (торсионными углами $\phi_i, \psi_i, \omega_i$). Вычисление этих углов позволяет получить ротационно-инвариантное описание формы цепи (оно не меняется при вращении или переносе белка в пространстве), что критически важно для нейросетевого анализа.

Конформация остова белковой цепи $-[N - C_\alpha - C]_i-$ практически полностью задается тремя двугранными (торсионными) углами $\phi_i, \psi_i, \omega_i$ для каждого остатка $i$.
Для их расчета используется ортогонализация векторов связей по методу Грама-Шмидта. Для четырех последовательно соединенных атомов $A-B-C-D$ с координатами положений $\mathbf{r}_A$, $\mathbf{r}_B$, $\mathbf{r}_C$, $\mathbf{r}_D$ определим направляющие векторы:
\begin{align}
    \mathbf{b}_0 &= \mathbf{r}_A - \mathbf{r}_B \\
    \mathbf{b}_1 &= \mathbf{r}_C - \mathbf{r}_B \\
    \mathbf{b}_2 &= \mathbf{r}_D - \mathbf{r}_C
\end{align}
Единичный направляющий вектор центральной связи запишется как $\mathbf{n}_1 = \mathbf{b}_1 / \|\mathbf{b}_1\|_2$. Спроецируем крайние векторы $\mathbf{b}_0$ и $\mathbf{b}_2$ на ортогональную $\mathbf{n}_1$ плоскость:
\begin{align}
    \mathbf{v} &= \mathbf{b}_0 - (\mathbf{b}_0 \cdot \mathbf{n}_1)\mathbf{n}_1 \\
    \mathbf{w} &= \mathbf{b}_2 - (\mathbf{b}_2 \cdot \mathbf{n}_1)\mathbf{n}_1
\end{align}
Значение двугранного угла $\theta(A, B, C, D) \in [-180^\circ, 180^\circ]$ вычисляется в градусах через функцию $\operatorname{atan2}$ (с переводом радианов в градусы):
\begin{equation}
    \theta(A, B, C, D) = \operatorname{atan2}\left( (\mathbf{n}_1 \times \mathbf{v}) \cdot \mathbf{w}, \mathbf{v} \cdot \mathbf{w} \right) \cdot \frac{180^\circ}{\pi}
\end{equation}
Двугранные углы остова вычисляются по формулам:
\begin{align}
    \phi_i &= \theta(\mathbf{r}_{C, i-1}, \mathbf{r}_{N, i}, \mathbf{r}_{C_\alpha, i}, \mathbf{r}_{C, i}) \\
    \psi_i &= \theta(\mathbf{r}_{N, i}, \mathbf{r}_{C_\alpha, i}, \mathbf{r}_{C, i}, \mathbf{r}_{N, i+1}) \\
    \omega_i &= \theta(\mathbf{r}_{C_\alpha, i-1}, \mathbf{r}_{C, i-1}, \mathbf{r}_{N, i}, \mathbf{r}_{C_\alpha, i})
\end{align}

\subsection{Карта Рамачандрана и выбросы геометрии}
Атомы имеют физический объем и не могут перекрываться в пространстве. Из-за этого вращение вокруг связей ограничено: аминокислоты могут принимать только определенные «разрешенные» комбинации углов $(\phi, \psi)$. Карта Рамачандрана — это распределение таких разрешенных зон (характерных для спиралей и листов). Вычисление «выбросов» (outliers) позволяет выявить участки структуры с физически невозможной или сильно напряженной геометрией, что указывает на ошибки генерации.

Стерическое отталкивание ограничивает возможные конформации углов $\phi, \psi$. Модель проверяет принадлежность углов трем классическим разрешенным структурным зонам:
\begin{itemize}
    \item \textbf{Зона $\alpha$-спиралей}: $\phi \in (-90^\circ, -30^\circ)$ и $\psi \in (-80^\circ, -10^\circ)$.
    \item \textbf{Зона $\beta$-листов}: $\phi \in (-170^\circ, -50^\circ)$ и $\psi \in (80^\circ, 180^\circ)$.
    \item \textbf{Зона левозакрученных спиралей}: $\phi \in (30^\circ, 90^\circ)$ и $\psi \in (20^\circ, 80^\circ)$.
\end{itemize}
Остаток помечается как выброс Рамачандрана (outlier), если пара $(\phi_i, \psi_i)$ выпадает за рамки всех разрешенных зон:
\begin{equation}
    \text{Outlier}_i = \neg \left( \text{in\_alpha}_i \lor \text{in\_beta}_i \lor \text{in\_lalpha}_i \right) \land \text{phi\_mask}_i \land \text{psi\_mask}_i \land M_i
\end{equation}
где $M_i$ — маска валидности остатка, а $\text{phi\_mask}_i$ и $\text{psi\_mask}_i$ — маски крайних остатков ($i > 0$ и $i < N-1$), исключающие ложные выбросы на концах цепи из-за паддинга. Выбросы свидетельствуют об ошибках в упаковке структуры.

\subsection{Восстановление координат виртуальных атомов}
При проектировании белков часто моделируют только центральный «скелет» (остов) белка без боковых радикалов, чтобы сэкономить вычислительные ресурсы. Однако для оценки стабильности структуры нужно знать: будут ли сталкиваться боковые цепи аминокислот (клэши) и могут ли образовываться водородные связи для удержания формы. Мы математически восстанавливаем положения ключевых атомов — виртуального $C_\beta$ (начало боковой цепи) и виртуального карбонильного $O$ (акцептор водородной связи) — для быстрой оценки этих параметров без построения полноатомной модели.

С целью анализа взаимодействия боковых цепей и водородного каркаса без явного дизайна полноатомной структуры, рассчитываются положения виртуальных атомов $C_\beta$ и $O$.

\paragraph{Виртуальный бета-углерод $C_\beta$.}
Положение атома $C_\beta$ восстанавливается на основе направлений смежных связей $\mathbf{v}_{N} = \mathbf{r}_{N, i} - \mathbf{r}_{C_\alpha, i}$ и $\mathbf{v}_{C} = \mathbf{r}_{C, i} - \mathbf{r}_{C_\alpha, i}$:
\begin{align}
    \mathbf{u}_{bis} &= \operatorname{normalize}\left(\operatorname{normalize}(\mathbf{v}_N) + \operatorname{normalize}(\mathbf{v}_C)\right) \\
    \mathbf{u}_{norm} &= \operatorname{normalize}\left(\operatorname{normalize}(\mathbf{v}_N) \times \operatorname{normalize}(\mathbf{v}_C)\right) \\
    \mathbf{d}_{C_\beta} &= -0.58 \cdot \mathbf{u}_{bis} + 0.81 \cdot \mathbf{u}_{norm} \\
    \mathbf{r}_{C_\beta, i} &= \mathbf{r}_{C_\alpha, i} + 1.53 \text{ \AA} \cdot \operatorname{normalize}(\mathbf{d}_{C_\beta})
\end{align}
Расстояние между виртуальными атомами $C_\beta$ менее $3.5\,\text{\AA}$ указывает на пространственный стерический конфликт (clash). Пары остатков с $|i-j| < 3$ из подсчета исключаются, поскольку их пространственная близость обусловлена геометрией самой цепи, а не конфликтом упаковки.

\paragraph{Виртуальный карбонильный кислород $O$.}
Карбонильный углерод находится в тригональной плоскости с связями на $C_\alpha$, $N_{i+1}$ и $O$ (углы $\approx 116^\circ$--$122^\circ$), поэтому направление на кислород лежит напротив биссектрисы направлений $\mathbf{u}_{CO} = \mathbf{r}_{C_\alpha, i} - \mathbf{r}_{C, i}$ и $\mathbf{v}_{CN} = \mathbf{r}_{N, i+1} - \mathbf{r}_{C, i}$:
\begin{align}
    \mathbf{d}_{O} &= \operatorname{normalize}\left(\operatorname{normalize}(\mathbf{u}_{CO}) + \operatorname{normalize}(\mathbf{v}_{CN})\right) \\
    \mathbf{r}_{O, i} &= \mathbf{r}_{C, i} - 1.23 \text{ \AA} \cdot \mathbf{d}_{O}
\end{align}
(знак <<минус>> принципиален: противоположное направление помещало бы кислород внутрь цепи между $C_\alpha$ и $N_{i+1}$). Для C-концевого остатка направление $\mathbf{v}_{CN}$ не определено и подставляется $\mathbf{u}_{CO}$; такой остаток не участвует в подсчете водородных связей, поскольку для него не существует пар с $j - i \ge 3$.
Водородная связь детектируется, если расстояние между $O_i$ и $N_j$ менее $3.5\,\text{\AA}$ (для $j-i \ge 3$).

\subsection{Прокси-дескрипторы плотности и упаковки}
Нативные белки сворачиваются в плотные шарики (глобулы), внутри которых гидрофобные остатки плотно упакованы в ядро, а гидрофильные выставлены наружу. Искусственные генеративные модели часто выдают «рыхлые» или полые структуры. Для отслеживания этого вычисляются: \textit{плотность контактов} (сколько соседей окружают остаток), \textit{экспонированность ядра} (насколько глубоко остаток погружен внутрь сферы), \textit{изгиб цепи} (кривизна «скелета») и \textit{фрагментная схожесть} (сравнение формы локальных окон из 9 остатков с реальными структурами из PDB). Это позволяет численно оценить плотность и корректность упаковки глобулы.

\begin{enumerate}
    \item \textbf{Плотность упаковки ядра (Contact Density)}: число атомов $C_\alpha$ других остатков на расстоянии менее $6, 8, 10\,\text{\AA}$:
    \begin{equation}
        D_{r, i} = \sum_{j: |i-j| > 3} [\|\mathbf{r}_{C_\alpha, i} - \mathbf{r}_{C_\alpha, j}\|_2 < r] \cdot M_j
    \end{equation}
    Мы считаем, сколько соседей окружают каждый остаток в сферах разного радиуса, исключая ближайших соседей по цепочке ($|i-j| > 3$), чтобы измерять именно пространственную упаковку, а не просто связность. У хорошо свернутого белка ядро упаковано очень плотно (много соседей), а поверхность — разреженно. Аномально низкая плотность внутри ядра указывает на «рыхлую» и нестабильную генерацию с пустотами.

    \item \textbf{Относительная экспонированность ядра (Relative Burial)}: расстояние от атома $C_{\alpha, i}$ до геометрического центра белка $\mathbf{r}_{center}$, нормированное на ожидаемый радиус глобулы $R_{expected} = 2.2 \cdot N_{eff}^{0.333} + 10^{-5}$:
    \begin{equation}
        B_i = \frac{\|\mathbf{r}_{C_\alpha, i} - \mathbf{r}_{center}\|_2}{R_{expected}}, \quad \mathbf{r}_{center} = \frac{1}{N_{eff}}\sum_k \mathbf{r}_{C_\alpha, k} \cdot M_k
    \end{equation}
    Показывает, насколько глубоко остаток погружен внутрь белкового шара. Значение около $0$ означает нахождение в самом центре ядра, а значение около $1$ или больше — на поверхности. Это позволяет модели проверять «эффект гидрофобного ядра» (правильно ли «водоотталкивающие» звенья спрятаны внутри структуры). Радиус нормируется на ожидаемый размер реального белка аналогичной длины ($R_{expected}$).

    \item \textbf{Локальный изгиб цепи (Local Bending)}: прокси-кривизна, выражаемая через расстояние с шагом 4:
    \begin{equation}
        K_i = \|\mathbf{r}_{C_\alpha, i+2} - \mathbf{r}_{C_\alpha, i-2}\|_2
    \end{equation}
    Мы измеряем расстояние по прямой между остатками через два шага назад и два шага вперед. Если это расстояние близко к максимальному ($\approx 15\,\text{\AA}$), то белковая цепь в этом месте идет прямо. Если расстояние маленькое, значит, цепь делает резкий разворот или петлю. Этот простой геометрический показатель заменяет вычисление кривизны и помогает выявлять неестественные изломы «скелета» белка.

    \item \textbf{Фрагментная схожесть (PDB-Fragment Similarity)}: 16-мерный дескриптор, получаемый линейным проецированием вектора попарных расстояний внутри локального окна размером $K=9$ остатков.

    Для каждого остатка $i \in \{1, \dots, N\}$ выделяется скользящее окно из 9 последовательных остатков, центрированное относительно $i$. Для обработки краев последовательности применяется дублирующий паддинг (replicate padding) на 4 элемента слева и справа. Пусть координаты атомов $C_\alpha$ в окне для остатка $i$ равны $\{\mathbf{x}_{i, 1}, \mathbf{x}_{i, 2}, \dots, \mathbf{x}_{i, 9}\}$, где $\mathbf{x}_{i, a} \in \mathbb{R}^3$.

    Для каждого такого окна рассчитывается матрица попарных расстояний $\mathbf{D}^{(i)} \in \mathbb{R}^{9 \times 9}$:
    \begin{equation}
        D^{(i)}_{ab} = \|\mathbf{x}_{i, a} - \mathbf{x}_{i, b}\|_2, \quad a, b \in \{1, \dots, 9\}
    \end{equation}
    Из полученной симметричной матрицы извлекаются элементы строго над главной диагональю, формирующие 36-мерный вектор попарных расстояний остатков в окне $\mathbf{d}_{pairs, i} \in \mathbb{R}^{36}$:
    \begin{equation}
        \mathbf{d}_{pairs, i} = \left[ D^{(i)}_{ab} \right]_{1 \le a < b \le 9}
    \end{equation}
    Для сжатия признаков и снижения шума вектор $\mathbf{d}_{pairs, i}$ линейно проецируется в 16-мерное пространство признаков $\mathbf{P}_{frag, i} \in \mathbb{R}^{16}$ с помощью PCA-матрицы проекции $\mathbf{W}_{PCA} \in \mathbb{R}^{16 \times 36}$ (что соответствует слою \texttt{Linear(36, 16, bias=False)} в классе \texttt{FoldabilityProxies} проекта):
    \begin{equation}
        \mathbf{P}_{frag, i} = \mathbf{W}_{PCA} \cdot \mathbf{d}_{pairs, i}
    \end{equation}
    Матрица $\mathbf{W}_{PCA}$ рассчитывается заранее методом главных компонент (PCA) на основе наборов структур реальных белков из базы PDB и фиксируется (замораживается) в процессе обучения основной модели.

    36 расстояний полностью определяют жесткую геометрию 9-остаточного фрагмента (его пространственную укладку). Линейная проекция сжимает эту информацию до 16 главных признаков, кодирующих, насколько локальная форма похожа на стандартные природные элементы вторичной структуры ($\alpha$-спирали, $\beta$-листы, петельные переходы).


\end{enumerate}

\subsection{Попарные признаки и построение графа контактов}
Белок — это не просто набор независимых точек, его структура держится на взаимодействиях между удаленными по цепи, но близкими в пространстве остатками. Мы вычисляем попарные характеристики взаимодействия (плавные RBF-расстояния, флаги непосредственного контакта, удаленность по последовательности). Чтобы нейросети не приходилось обрабатывать взаимодействия каждой пары остатков с каждой (что требует квадратичного времени $O(N^2)$ и неэффективно для длинных белков), мы строим граф ближайших соседей (kNN-граф, $k=16$). В нем ребра соединяют только пространственно близкие остатки, что позволяет свести сложность анализа к линейной $O(N)$.

Помимо узловых признаков, рассчитываются характеристики взаимодействия пар остатков $(i, j)$:
\begin{itemize}
    \item \textbf{RBF-разложение расстояния (Radial Basis Function Expansion)}:
    Расстояние $d_{ij} = \|\mathbf{r}_{C_\alpha, i} - \mathbf{r}_{C_\alpha, j}\|_2$ между альфа-углеродами остатков $i$ и $j$ кодируется вектором из $B=16$ значений с помощью радиально-базисных функций, равномерно распределенных на интервале от $d_{min} = 2.0\,\text{\AA}$ до $d_{max} = 22.0\,\text{\AA}$:
    \begin{equation}
        \mathrm{RBF}_{ij, k} = \exp\left(-\gamma (d_{ij} - \mu_k)^2\right), \quad k \in \{1, \dots, B\}
    \end{equation}
    где центры функций $\mu_k$ задаются:
    \begin{equation}
        \mu_k = d_{min} + (k - 1) \cdot \frac{d_{max} - d_{min}}{B - 1}
    \end{equation}
    а коэффициент ширины ядра $\gamma$ определяется нормировочным шагом $\Delta = \frac{d_{max} - d_{min}}{B} = 1.25\,\text{\AA}$:
    \begin{equation}
        \gamma = \frac{1}{\Delta^2} = \frac{1}{(1.25\,\text{\AA})^2} = 0.64\,\text{\AA}^{-2}
    \end{equation}

    \textit{Справка по RBF-разложению:} В задачах трехмерного анализа молекулярных структур прямое использование скалярного расстояния $d_{ij}$ в качестве входного признака малоэффективно, так как нейронные сети плохо аппроксимируют высокочастотные локальные изменения расстояний с помощью одного вещественного числа. RBF-разложение (или «плавное бинирование», soft binning) представляет собой отображение расстояния в высокомерный вектор признаков. Вместо жесткого разбиения на интервалы, где переход через границу ячейки разрывен, RBF обеспечивает гладкое, дифференцируемое представление: при изменении расстояния плавно меняются активации сразу нескольких соседних гауссовых каналов.
    \item \textbf{Бинарный контактный флаг}: $[d_{ij} < 8.0\,\text{\AA}]$.
    \item \textbf{Признаки взаимного положения по последовательности}: относительное расстояние по цепи $|i-j| / 32$, индикатор соседства $|i-j|=1$ и индикатор принадлежности локальному окну $|i-j| \le 4$.
\end{itemize}
На основе попарных расстояний строится разреженный kNN-граф контактов ($k=16$ ближайших соседей по расстоянию $C_\alpha - C_\alpha$ для каждого остатка), что позволяет свести квадратичную сложность попарного анализа к линейной по числу остатков.

\newpage

% --- КОНСТРУИРОВАНИЕ ДАТАСЕТА И НЕГАТИВНОЕ СЭМПЛИРОВАНИЕ ---
\section{Конструирование датасета и негативное сэмплирование}

\subsection{Сбор нативных структур белков}
Сбор исходных нативных структур (положительный класс, \texttt{label} $= 1.0$) производился с помощью RCSB API. Для обеспечения качества исходных данных выборка фильтровалась по следующим критериям:
\begin{itemize}
    \item Метод определения: рентгеноструктурный анализ (X-ray diffraction) высокого разрешения ($\le 2.0\,\text{\AA}$).
    \item Отсутствие длинных петель с пропущенными координатами атомов остова (проверялось ограничением расстояния ковалентной связи $C_i - N_{i+1} \le 2.0\,\text{\AA}$).
    \item Диапазон длин белков от $50$ до $700$ остатков.
\end{itemize}
Координаты атомов остова динамически центрируются относительно начала координат непосредственно в загрузчике данных (DataLoader) перед подачей в модель.

\subsection{Многоуровневая генерация декоев}
Для эффективного обучения модели детекции дефектов упаковки, на основе каждого нативного белка было сгенерировано по два декоя с использованием 7 деформационных стратегий, имитирующих реальные ошибки генеративного дизайна:
\begin{enumerate}
    \item \textbf{easy\_global\_noise} (простой декой, класс 1): наложение гауссова шума со среднеквадратичным (стандартным) отклонением $1.5\,\text{\AA}$ на все атомы остова белка.
    \item \textbf{easy\_chain\_break} (простой декой, класс 1): разрезание цепи в случайном месте со смещением одной из частей на $20.8$--$31.2\,\text{\AA}$ (путём независимого сдвига по каждой координатной оси на $12$--$18\,\text{\AA}$).
    \item \textbf{hard\_core\_unpacked} (сложный декой, класс 2): центроидное расширение (радиальное раздутие) цепи на 10--25\%, разрушающее компактную упаковку гидрофобного ядра.
    \item \textbf{hard\_false\_compact} (сложный декой, класс 2): разделение остова на 2--4 случайных сегмента с поворотом блоков и повторной упаковкой в рамках плотной сферы малого радиуса (4.5--8.5\,\text{\AA}).
    \item \textbf{hard\_near\_native} (сложный декой, класс 3): анизотропное сжатие или растяжение остова вдоль случайных осей (на 2--5\% по двум осям и на 0--2\% по третьей оси) с произвольными поворотами на угол до $180^\circ$.
    \item \textbf{borderline\_hinge\_defect} (пограничный декой, класс 4): поворот хвоста структуры относительно случайной оси (шарнира) на угол от $8^\circ$ до $22^\circ$.
    \item \textbf{borderline\_local\_fragment\_rotation} (пограничный декой, класс 4): локальный поворот короткого фрагмента (3--8 остатков) на угол $6^\circ$--$18^\circ$ вокруг продольной оси.
\end{enumerate}

\subsection{Предотвращение утечки данных (Data Leakage)}
С целью исключения переобучения модели, распределение сэмплов по сплитам (70\% обучающая, 15\% валидационная, 15\% тестовая выборки) выполнялось строго сгруппированно по идентификатору родительской цепи белка (\texttt{group\_id} на основе сочетания PDB ID и Chain ID). Все декои, сгенерированные из конкретной нативной цепи белка, и сама эта цепь попадали строго в одну выборку. Случайный сплит привел бы к ситуации, когда модель видит схожие искажения одной и той же цепи как в обучающей, так и в тестовой выборках, что исказило бы оценку точности модели.

\newpage

% --- АЛГОРИТМЫ РАСЧЕТА ДЕСКРИПТОРОВ ---
\section{Алгоритмы расчета дескрипторов}

В данном разделе приводится структурированный псевдокод алгоритмов извлечения геометрических и биофизических признаков белков de novo.

\subsection{Алгоритм экстракции геометрических признаков}
Алгоритм вычисляет торсионные углы, выбросы Рамачандрана, плотность клэшей и количество водородных связей. Соответствует классу \texttt{Backbone}\-\texttt{Geometry}\-\texttt{Extractor}.

\begin{algorithm}[h]
\caption{Извлечение геометрии остова белка}
\label{alg:geometry}
\begin{algorithmic}[1]
\Statex \textbf{Вход:} Координаты остова $C \in \mathbb{R}^{B \times N \times 3 \times 3}$, Маска валидности $M \in \mathbb{B}^{B \times N}$
\Statex \textbf{Выход:} Набор локальных геометрических признаков
\State Разделить координаты на атомы: $N_{at} \gets C[:, :, 0]$, $C_\alpha \gets C[:, :, 1]$, $C_{carb} \gets C[:, :, 2]$
\State Вычислить сырые торсионные углы $\phi_{raw}, \psi_{raw}, \omega_{raw}$ по формуле двугранного угла (метод Грама-Шмидта):
\State \quad $\phi_{raw} \gets \operatorname{dihedral}(C_{carb}[:, :-1], N_{at}[:, 1:], C_\alpha[:, 1:], C_{carb}[:, 1:])$
\State \quad $\psi_{raw} \gets \operatorname{dihedral}(N_{at}[:, :-1], C_\alpha[:, :-1], C_{carb}[:, :-1], N_{at}[:, 1:])$
\State \quad $\omega_{raw} \gets \operatorname{dihedral}(C_\alpha[:, :-1], C_{carb}[:, :-1], N_{at}[:, 1:], C_\alpha[:, 1:])$
\State Выполнить паддинг и маскирование для выравнивания по размерности $N$:
\State \quad $\phi \gets \operatorname{pad\_left}(\phi_{raw}, 1)$, $\psi \gets \operatorname{pad\_right}(\psi_{raw}, 1)$, $\omega \gets \operatorname{pad\_left}(\omega_{raw}, 1)$
\State \quad $\text{phi\_mask}_i \gets [i > 0]$, $\text{psi\_mask}_i \gets [i < N-1]$ \Comment{маски исключения крайних остатков}
\State Вычислить признаки выбросов Рамачандрана:
\State \quad $\text{in\_alpha}_i \gets (\phi_i \in (-90^\circ, -30^\circ)) \land (\psi_i \in (-80^\circ, -10^\circ))$
\State \quad $\text{in\_beta}_i \gets (\phi_i \in (-170^\circ, -50^\circ)) \land (\psi_i \in (80^\circ, 180^\circ))$
\State \quad $\text{in\_lalpha}_i \gets (\phi_i \in (30^\circ, 90^\circ)) \land (\psi_i \in (20^\circ, 80^\circ))$
\State \quad $\text{outliers}_i \gets \neg(\text{in\_alpha}_i \lor \text{in\_beta}_i \lor \text{in\_lalpha}_i) \land \text{phi\_mask}_i \land \text{psi\_mask}_i \land M_i$
\State Восстановить координаты виртуального $C_\beta$:
\State \quad $\mathbf{v}_N \gets N_{at} - C_\alpha$, $\mathbf{v}_C \gets C_{carb} - C_\alpha$
\State \quad $\mathbf{u}_{bis} \gets \operatorname{normalize}\left(\operatorname{normalize}(\mathbf{v}_N) + \operatorname{normalize}(\mathbf{v}_C)\right)$
\State \quad $\mathbf{u}_{norm} \gets \operatorname{normalize}\left(\operatorname{normalize}(\mathbf{v}_N) \times \operatorname{normalize}(\mathbf{v}_C)\right)$
\State \quad $\mathbf{d}_{C_\beta} \gets -0.58 \cdot \mathbf{u}_{bis} + 0.81 \cdot \mathbf{u}_{norm}$
\State \quad $C_\beta \gets C_\alpha + 1.53\,\text{\AA} \cdot \operatorname{normalize}(\mathbf{d}_{C_\beta})$
\State Восстановить координаты виртуального кислорода $O$:
\State \quad $\mathbf{u}_{CO} \gets C_\alpha - C_{carb}$, $\mathbf{v}_{CN} \gets \operatorname{roll}(N_{at}, \text{shift}=-1) - C_{carb}$
\State \quad $\mathbf{d}_O \gets \operatorname{normalize}\left(\operatorname{normalize}(\mathbf{u}_{CO}) + \operatorname{normalize}(\mathbf{v}_{CN})\right)$
\State \quad $O \gets C_{carb} + 1.23\,\text{\AA} \cdot \mathbf{d}_O$
\State Вычислить водородные связи:
\State \quad $D_{ON} \gets \operatorname{cdist}(O, N_{at})$
\State \quad $\text{hb\_mask}_{ij} \gets (D_{ON, ij} < 3.5\,\text{\AA}) \land (j - i \ge 3) \land M_i \land M_j$
\State \quad $\text{hbond\_count}_i \gets \sum_{j} \text{hb\_mask}_{ij}$
\State Вычислить стерические клэши:
\State \quad $D_{C_\beta} \gets \operatorname{cdist}(C_\beta, C_\beta)$
\State \quad $\text{clash\_mask}_{ij} \gets (D_{C_\beta, ij} < 3.5\,\text{\AA}) \land (i \neq j) \land M_i \land M_j$
\State \quad $\text{clash\_count}_i \gets \sum_{j} \text{clash\_mask}_{ij}$
\State \Return $\{\phi, \psi, \omega, \text{outliers}, \text{hbond\_count}, \text{clash\_count}\}$
\end{algorithmic}
\end{algorithm}

\newpage

\subsection{Алгоритм экстракции признаков свертываемости}
Алгоритм извлекает признаки плотности упаковки ядра, экспонированности растворителю и фрагментной схожести с использованием PCA-проекции. Соответствует классу \texttt{Foldability}\-\texttt{Proxies}.

\begin{algorithm}[h]
\caption{Извлечение признаков свертываемости}
\label{alg:foldability}
\begin{algorithmic}[1]
\Statex \textbf{Вход:} Координаты $C_\alpha \in \mathbb{R}^{B \times N \times 3}$, Маска $M \in \mathbb{B}^{B \times N}$, Матрица расстояний $D \in \mathbb{R}^{B \times N \times N}$
\Statex \textbf{Выход:} Прокси-признаки свертываемости
\State Инициализировать маску дальних контактов: $V_{ij} \gets (|i - j| > 3) \land M_i \land M_j$
\State Вычислить плотность контактов:
\For{каждого радиуса $r \in \{6, 8, 10\}$}
    \State $\text{density}_{r, i} \gets \sum_{j=1}^N [D_{ij} < r] \cdot V_{ij}$
\EndFor
\State Оценить относительную экспонированность:
\State \quad $\mathbf{r}_{center} \gets \left(\sum_j C_{\alpha, j} \cdot M_j\right) / \left(\sum_j M_j + 10^{-8}\right)$
\State \quad $d_{center, i} \gets \|C_{\alpha, i} - \mathbf{r}_{center}\|_2$
\State \quad $R_{expected} \gets 2.2 \cdot \left(\sum_j M_j\right)^{0.333} + 10^{-5}$
\State \quad $\text{burial}_i \gets d_{center, i} / R_{expected}$
\State Вычислить фрагментную схожесть:
\State \quad Выделить скользящие окна размера $K=9$ по $C_\alpha$ с replicate padding $(4, 4)$
\State \quad Для каждого окна собрать вектор расстояний $\mathbf{d}_{pairs} \in \mathbb{R}^{36}$ (верхний треугольник)
\State \quad Спроецировать: $P_{frag} \gets \mathbf{W}_{PCA} \cdot \mathbf{d}_{pairs}$ (где $\mathbf{W}_{PCA} \in \mathbb{R}^{16 \times 36}$)
\State Вычислить локальную кривизну (local bending):
\State \quad $K_i \gets \|C_{\alpha, i+2} - C_{\alpha, i-2}\|_2$ (с дополнением нулями для краев)
\State \Return $\{\text{density}_{6,8,10}, \text{burial}, P_{frag}, K\}$
\end{algorithmic}
\end{algorithm}

\newpage

\subsection{Алгоритм построения парных признаков и kNN-графа}
Алгоритм вычисляет RBF-разложение расстояний между остатками и строит список ребер и их атрибутов для $k$-ближайших соседей. Соответствует классу \texttt{Pair}\-\texttt{Feature}\-\texttt{Builder}.

\begin{algorithm}[h]
\caption{Построение парных признаков и kNN-графа}
\label{alg:pair_features}
\begin{algorithmic}[1]
\Statex \textbf{Вход:} Координаты $C_\alpha \in \mathbb{R}^{B \times N \times 3}$, Маска $M \in \mathbb{B}^{B \times N}$, Матрица расстояний $D \in \mathbb{R}^{B \times N \times N}$
\Statex \textbf{Выход:} Признаки пар $E_{attr}$, список ребер графа $E_{idx}$
\State Вычислить RBF-разложение для всех пар расстояний:
\State \quad $RBF_{ij} \gets \exp\left(-\gamma \cdot (D_{ij} - \vec{\mu})^2\right) \in \mathbb{R}^{16}$
\State Сформировать признаки пар:
\State \quad $C_{flag, ij} \gets [D_{ij} < 8.0]$ \Comment{флаг контакта}
\State \quad $S_{sep, ij} \gets \min(1.0, |i-j| / 32.0)$ \Comment{нормализованное расстояние по цепи}
\State \quad $S_{adj, ij} \gets [|i-j| == 1]$ \Comment{индикатор смежности}
\State \quad $W_{local, ij} \gets [|i-j| \le 4]$ \Comment{индикатор локального окна}
\State \quad $P_{ij} \gets \operatorname{Concat}(RBF_{ij}, C_{flag, ij}, S_{sep, ij}, S_{adj, ij}, W_{local, ij}) \in \mathbb{R}^{20}$
\State Инициализировать списки ребер графа $E_{idx} \gets []$ и атрибутов ребер $E_{attr} \gets []$
\For{каждого белка $b$ в батче $B$}
    \State Найти валидные остатки по маске $M_b$ (всего $N_b$ узлов)
    \State Создать матрицу расстояний для валидных остатков, заменяя диагональ и паддинги на $+\infty$:
    \State \quad $D'_{ij} \gets D_{b, ij}$ с маскированием по $M_{b, i} \land M_{b, j}$ и $i \neq j$
    \State Для каждого узла отобрать $k=16$ ближайших соседей по матрице $D'$:
    \State \quad Выбрать ребра к $k$ соседям с конечными расстояниями
    \State Сформировать локальный список ребер $E_{idx, b}$ и собрать атрибуты ребер $E_{attr, b}$ из $P_b$
    \State Добавить $E_{idx, b}$ и $E_{attr, b}$ в глобальные списки
\EndFor
\State \Return $\{E_{idx}, E_{attr}\}$
\end{algorithmic}
\end{algorithm}

\newpage

% --- ЗАКЛЮЧЕНИЕ ---
\section{Заключение}
В ходе выполнения данной курсовой работы были спроектированы и программно реализованы методы биофизического препроцессинга пространственных структур остова белков de novo и конвейер генерации сбалансированных наборов данных.

Основные результаты работы:
\begin{enumerate}
    \item Разработан модуль биофизического анализа геометрии белкового остова, вычисляющий $31$ инвариантный признак узлов и $20$ признаков пар, отражающих локальные углы, стерическое напряжение и контакты.
    \item Реализован устойчивый алгоритм восстановления координат виртуальных атомов $C_\beta$ и $O$ для быстрой оценки водородного каркаса и стерической плотности без построения боковых цепей.
    \item Предложены 7 стратегий негативного сэмплирования, моделирующие различные дефекты упаковки и связности остова, что позволило получить сбалансированный датасет декоев.
    \item Разработана и обоснована методика группового разбиения данных по родительским цепям белков (PDB ID + Chain ID), предотвращающая утечку данных на этапе обучения.
\end{enumerate}

Разработанные алгоритмы извлечения признаков и подготовки данных представляют самостоятельную ценность для систем препроцессинга геометрических нейросетей в структурной биоинформатике.

\newpage

% --- СПИСОК ЛИТЕРАТУРЫ ---
\begin{thebibliography}{99}
\addcontentsline{toc}{section}{Список литературы}

\bibitem{alphafold}
Jumper, J., Evans, R., Pritzel, A. et al. Highly accurate protein structure prediction with AlphaFold. \textit{Nature} 596, 583–589 (2021).

\bibitem{rfdiffusion}
Watson, J. L., Juergens, D., Bennett, N. R. et al. De novo design of protein structure and function with RFdiffusion. \textit{Nature} 620, 1089–1100 (2023).

\bibitem{chroma}
Ingraham, J., Barzilay, R., Regina, A. et al. Illuminating protein space with a programmable generative model. \textit{bioRxiv} 2022.12.01.518682 (2022).

\bibitem{biotite}
Kunzmann, P., Hamacher, K. Biotite: a unifying framework for biophysical data in Python. \textit{BMC Bioinformatics} 19, 1–10 (2018).

\end{thebibliography}

\end{document}
